AI-generated image detection has advanced to the point where, on the standard GenImage benchmark, recent detectors report near-saturated cross-generator accuracy; our own deliberately conventional, single-GPU detector reaches 99.67\% mean accuracy on held-out GenImage generators. A saturated benchmark, however, leaves open what such numbers are made of: a high score can reflect either the detection of generative traces or the exploitation of incidental dataset artifacts. We study this question from inside the model. Our detector---a dual-stream backbone with a six-channel concept bottleneck and a linear head---has one essential property: the prediction passes through six separately-testable evidence channels and nothing else. We use this bottleneck not as an interpretability device but as an audit instrument, measuring each channel's contribution causally, by ablation and counterfactual, rather than reading it from weights or post-hoc narration. The audit shows that the detector's decision turns on a frequency/compression cue pair (jpeg-quant and freq-radial, $r{=}{-}0.80$)---localized by intervention; the carrier label shifts across seeds, but the pair-level reliance is stable, a finding weight inspection alone would miss---and that this reliance persists even after Grommelt-style JPEG debiasing neutralizes the obvious container cue (replicated across three seeds, with seed-dependent strength), so the compression axis is not reducible to the image-format label---though we cannot exclude a residual JPEG re-encoding signature as the surviving cue. The headline accuracy is thus partly confound-inflated: it drops substantially on a container-mismatched benchmark, though multi-generator training still buys genuine partial cross-architecture generalization. We extend prior dataset-level analyses of GenImage confounds to a model-internal account---where the reliance lives, how large it is, and whether it survives debiasing---and argue that a concept bottleneck's value lies in being a testability surface. A no-bottleneck baseline matches this accuracy (99.60\% OOD) and is equally confound-susceptible, confirming the bottleneck adds audit visibility at zero accuracy cost. We report competitive, not state-of-the-art, detection.
